\section{The Eighth Day: The Octave of the Nativity, the Mother of God, and the Circumcision}\label{sec:jan-1}

Eight days after Bethlehem, the Christmas Octave reaches its summit. Under the
1960 \textit{Codex Rubricarum} (\textit{Variationes} nn.~6 and~12; General
Rubric~91, row~5), 1~January bears the restored ancient title \textbf{\textit{In
Octava Nativitatis Domini}}, the Octave Day of the Nativity of the Lord, raised
to the \textbf{first class} in white vestments, with the station at Rome's
oldest Marian basilica across the Tiber, \textit{Sancta Maria trans Tiberim}.
No matter what day of the week 1~January falls upon---even when it is a
Sunday---it takes complete precedence with no commemorations.

\subsection{Against the Pagan Kalends: St.~Almachius and the Early Church}

Before 1~January ever looked calm or domestic, it was a battleground. Across
the Roman Empire, the Kalends of January opened the civil year with idolatrous
processions, divination, cross-dressing revels, and drunken excess. The ancient
Church answered not by ignoring the date, but by planting the Octave of the
Nativity directly across it.

Open the \textit{Roman Martyrology} to 1~January and an astonishing entry
stands immediately beneath the heading of the feast: St.~Almachius, Martyr,
slain by gladiators at Rome at the command of Alypius, Prefect of the City,
``because he said: `This is the Octave of the Lord's Birth, cease from your
idolatrous superstitions, and your unclean sacrifices.''' In fifth-century
Ravenna, St.~Peter Chrysologus opened Sermon~155 (\textit{On the Kalends of
January}) with blunt eyewitness alarm: wherever Christ is born for our
salvation, the devil raises monstrous spectacles to mock religion, turn sanctity
into sacrilege, and fashion an insult against God out of the honor due to him,
as pagans drag their idols through the streets in studied disgrace. In the Old
Gelasian Sacramentary, right beside the Mass \textit{In Octavas Domini}, stands
a second Mass headed simply \textit{Prohibendum ab Idolis}---``To keep away from
idols.'' And in~567 the Second Council of Tours ordered a three-day fast and the
Mass of the Circumcision to draw Christians away from New Year's carousing.
Every 1~January, when the Church celebrates the Octave Day in quiet white
splendor, she still answers the noise of the world with the peace of Bethlehem.

\subsection{Rome's Oldest Marian Feast and the Five Antiphons}

Alongside that early vigil against paganism, 1~January preserves Rome's oldest
solemnity of the Mother of God, \textit{Natale Sanctae Mariae}, kept at Santa
Maria in Trastevere before the seventh-century Byzantine Marian feasts reached
the West (Section~\ref{sec:history}). That ancient Roman stratum explains why
the 1962 Mass and Office of 1~January are saturated with the Divine Maternity.

The Mass opens with the Day Mass Introit \textit{Puer natus est}, and then
prays the ancient Marian Collect that also closes \textit{Alma Redemptoris
Mater} every night of Christmastide:

\begin{massbox}{Collect of the Octave Day of the Nativity (1 January)}
\textbf{Latin} (1962 Roman Missal). Deus, qui salutis aeternae, beatae Mariae
virginitate fecunda, humano generi praemia praestitisti: tribue, quaesumus; ut
ipsam pro nobis intercedere sentiamus, per quam meruimus auctorem vitae
suscipere, Dominum nostrum Jesum Christum, Filium tuum: Qui tecum vivit et
regnat in unitate.

\textbf{English} (Bute, 1908). O GOD, Who, by the fruitful virginity of the
Blessed Mary, hast given unto mankind the rewards of everlasting life; grant, we
beseech Thee, that we may continually feel the might of her intercession,
through whom we have worthily received the Author of our life, our Lord JESUS
Christ, Thy Son, Who liveth and reigneth with Thee, in the unity of the Holy
Ghost, one God, world without end. Amen.
\end{massbox}

The Postcommunion (\textit{Haec nos communio}) repeats the same Marian note,
praying: ``May this communion, O Lord, cleanse us from sin and, by the
intercession of the Blessed Virgin Mary, Mother of God, make us partakers of the
heavenly remedy.''

In the Divine Office, First Vespers, Lauds, and Second Vespers of 1~January
sing the five psalms of the Blessed Virgin with the five celebrated antiphons
of the wonderful exchange (printed in full in Appendix~\ref{app:prayers}). The
first proclaims the whole mystery of Christmas: \textit{O admirabile commercium:
Creator generis humani, animatum corpus sumens, de Virgine nasci dignatus est:
et procedens homo sine semine, largitus est nobis suam Deitatem}---``How
wonderful is the dealing here! The Maker of man hath taken to Himself a living
body, and hath been pleased to be born of the Virgin: He is come forth conceived
without seed, and hath made us partakers of His Divine nature.'' The next four
unfold every Old Testament figure of the Virgin Mother: the rain falling quietly
on Gideon's fleece (\textit{Quando natus es}, Ps~71~[72]:6); the burning bush
seen by Moses unconsumed (\textit{Rubum quem viderat Moyses}, Ex~3:2); the Root
of Jesse blossoming and the Star rising out of Jacob (\textit{Germinavit radix
Jesse}); and Mary bearing the Savior whom John points out as the Lamb of God
(\textit{Ecce Maria genuit}). At Matins, the Sixth Responsory
(\textit{Benedicta et venerabilis es, Virgo Maria}) marvels at the Child of whom
it sings, ``He Whose glory filled the heavens lay in a manger,'' and the Eighth
Responsory (\textit{Nesciens mater virgo virum}) contemplates Mary: ``She fed
the Saviour of the world, The King of Angel hosts above, Jesus, our Redeemer
blest, From the fountain of her breast.''

\subsection{The Circumcision and the First Drops of Blood}

Although the 1960 rubrics restored the feast's ancient Roman title, \textit{In
Octava Nativitatis Domini}, the Gospel of the day remains the single, dense
verse of St.~Luke on the eighth day (Luke~2:21):

\begin{quote}
And after eight days were accomplished, that the child should be circumcised,
his name was called JESUS, which was called by the angel before he was conceived
in the womb.
\end{quote}

Joined to the Midnight Epistle (Tit~2:11--15) and the Alleluia from Hebrews
(\textit{Multifariam olim Deus}, Heb~1:1--2), this single verse draws out some
of the richest theology in the patristic and medieval tradition.

Why did the sinless Lawgiver submit to the knife of the Old Law? In the
\textit{Summa Theologiae} (Tertia Pars, q.~37, a.~1), St.~Thomas Aquinas lists
seven reasons: to prove the reality of his human flesh against heresies that
called his body an illusion; to approve the God-given Law of Moses; to prove
his descent from Abraham; to take away every excuse for rejecting him; to teach
us obedience on the exact appointed eighth day; to refuse no remedy ordained for
sinful flesh; and by bearing the burden of the Law himself, to redeem those who
were under the Law (Gal~4:4--5).

In the Greek East, St.~Cyril of Alexandria uncovered a luminous Old Testament
type in Sermon~3 of his \textit{Commentary on Luke} (on Luke~2:21--24):

\begin{quote}
For on the eighth day Christ arose from the dead, and gave us the spiritual
circumcision... And of this again, that Jesus [Joshua] of old, who was captain
after Moses, was a type. For he first of all led the children of Israel across
the Jordan: and then having halted them, immediately circumcised them with
knives of stone. So when we have crossed the Jordan, Christ circumcises us with
the power of the Holy Ghost, not purifying the flesh, but rather cutting off the
defilement that is in our souls.
\end{quote}

St.~John of Damascus adds in \textit{De Fide Orthodoxa} (Book~IV, Chapter~25)
that just as bodily circumcision removes only a useless superfluity, so holy
Baptism cuts away sin---``the superfluous part of desire and not useful
desire.'' In the Roman Breviary's Seventh and Eighth Lessons for 1~January,
St.~Ambrose carries the eighth day all the way to the end of time: because our
body and mind remain inclined to sin in this life, ``the circumcision of the
eighth day is also a type of that complete cleansing from sin which we shall
have at the resurrection.'' Meanwhile in the Second Nocturn, Pope St.~Leo the
Great's Sermon~27 (section~1) guards the confession of Chalcedon: one and the
same Person is true God and true man, who suffers the wound of the flesh while
remaining sovereign Lord of heaven.

Preaching his Second Sermon for the Circumcision, St.~Bernard of Clairvaux saw
Calvary already shining in the eighth day: ``Very properly, therefore, did the
Child Who was born for us receive the name of Jesus, that is, Saviour, at the
time of His circumcision; because it was then He began the work of our
salvation, by pouring out the first instalment of His immaculate Blood.''

At Mass on 1~January, the proper \textit{Communicantes} (\textit{diem
sacratissimum}) is prayed for the last time in the season. As Pope Leo~XIII
urged in \textit{Tametsi} (1900) at the turning of a new century, the Christian
steps across the threshold of a new civil year under no blind fate, but under
the rule of Christ the Redeemer, whose Blood was first shed on the eighth day
and whose saving Name is proclaimed before the world.
